\documentclass[../BTL.tex]{subfiles}
\begin{document}
\chapter{Phiên bản công nghệ sử dụng}
\label{phu-luc:versions}

Phụ lục này tổng hợp phiên bản được khai báo trong mã nguồn, tệp khóa phụ thuộc và cấu hình triển khai của hệ thống. Với thư viện được quản lý qua BOM, bảng ghi phiên bản BOM thay vì tự suy ra phiên bản riêng của từng thư viện. Phiên bản container được ghi đúng theo image tag trong cấu hình; tag không chỉ rõ bản vá thì không suy diễn thêm. Với công cụ cài qua trình quản lý gói của hệ điều hành, bảng ghi phiên bản đang cài trên máy tại thời điểm viết báo cáo kèm cách xác minh, không chuyển thành yêu cầu phiên bản cho dự án.

\section{Ngôn ngữ và công nghệ backend}

\begin{table}[H]
\centering
\small
\begin{tabularx}{\textwidth}{|p{4.0cm}|p{2.0cm}|X|}
\hline
Công nghệ & Phiên bản & Phạm vi / nguồn khai báo \\ \hline
Java & 21 & Thuộc tính \textnormal{java.version} trong các tệp \textnormal{pom.xml}; image xây dựng và chạy dùng Eclipse Temurin 21. \\ \hline
Spring Boot & 4.1.0 & Phiên bản parent Maven dùng chung cho sáu dịch vụ backend. \\ \hline
Spring Cloud & 2025.1.2 & BOM quản lý các thư viện Spring Cloud trong sáu dịch vụ. \\ \hline
Spring AI & 2.0.2 & BOM của \textnormal{ai-assistant-service}. \\ \hline
Apache Maven & 3.9.16 & Maven Wrapper trong sáu dịch vụ backend. \\ \hline
MapStruct & 1.6.3 & Khai báo trong \textnormal{course-service}, \textnormal{executor-service}, \textnormal{result-service} và \textnormal{submission-service}. \\ \hline
Springdoc OpenAPI & 3.1.0 & Khai báo trong \textnormal{course-service}, \textnormal{executor-service}, \textnormal{result-service} và \textnormal{submission-service}. \\ \hline
Spring AI OpenAI starter & 2.0.2 & Phiên bản do BOM Spring AI quản lý trong \textnormal{ai-assistant-service}. \\ \hline
Testcontainers & Theo BOM Spring Boot & Thư viện kiểm thử tích hợp; phiên bản do quản lý phụ thuộc Spring Boot cung cấp. \\ \hline
AWS SDK for Java S3 & 2.25.15 & Thư viện client S3 dùng để gọi RustFS trong \textnormal{submission-service}. \\ \hline
MinIO Java SDK & 8.5.14 & Thư viện client S3 dùng để gọi RustFS trong \textnormal{submission-service} và \textnormal{executor-service}. \\ \hline
JsonPath & 2.9.0 & Thư viện đọc đường dẫn JSON trong \textnormal{executor-service}. \\ \hline
Micrometer datasource observation & 2.3.0 & Thư viện quan sát lệnh cơ sở dữ liệu được khai báo trong các dịch vụ. \\ \hline
\end{tabularx}
\caption{Phiên bản ngôn ngữ và công nghệ backend.}
\label{tab:versions-backend}
\end{table}

Các thư viện Spring Boot và Spring Cloud không ghi phiên bản riêng trong từng dependency được quản lý theo BOM tương ứng; vì vậy phiên bản BOM là mốc khai báo có thể xác minh trực tiếp trong cấu hình dự án.

\section{Giao diện web}

Các phiên bản đã khóa trong \textnormal{package-lock.json} được dùng khi có khác biệt với khoảng phiên bản cho phép trong \textnormal{package.json}.

\begin{table}[H]
\centering
\small
\begin{tabularx}{\textwidth}{|p{3.5cm}|p{2.4cm}|X|}
\hline
Công nghệ & Phiên bản & Ghi chú \\ \hline
React & 19.2.8 & Thư viện giao diện và kết xuất. \\ \hline
TypeScript & 6.0.3 & Ngôn ngữ và bộ kiểm tra kiểu khi xây dựng giao diện. \\ \hline
Vite & 8.2.1 & Công cụ phát triển và đóng gói giao diện. \\ \hline
Ant Design & 6.6.5 & Bộ thành phần giao diện. \\ \hline
React Router & 8.4.0 & Thư viện định tuyến phía trình duyệt. \\ \hline
TanStack Query & 5.104.1 & Quản lý và lưu đệm trạng thái dữ liệu từ API. \\ \hline
Keycloak JavaScript adapter & 26.2.4 & Thư viện tích hợp đăng nhập trên giao diện. \\ \hline
Vitest & 5.1.0 & Phiên bản khóa trong \textnormal{package-lock.json}; \textnormal{package.json} khai báo khoảng \textnormal{\textasciicircum 5.0.5}. \\ \hline
Axios & 1.20.0 & Gửi yêu cầu tới API; phiên bản khóa trong tệp khóa. \\ \hline
i18next & 26.4.2 & Thư viện đa ngôn ngữ; phiên bản khóa trong tệp khóa. \\ \hline
React i18next & 17.0.15 & Tích hợp đa ngôn ngữ với React; phiên bản khóa trong tệp khóa. \\ \hline
React Markdown & 10.1.0 & Hiển thị nội dung Markdown; phiên bản khóa trong tệp khóa. \\ \hline
Monaco Editor React & 4.7.0 & Tích hợp trình soạn thảo Monaco; phiên bản khóa trong tệp khóa. \\ \hline
Ant Design Icons & 6.3.4 & Bộ biểu tượng giao diện; phiên bản khóa trong tệp khóa. \\ \hline
\end{tabularx}
\caption{Phiên bản các công nghệ chính của giao diện web.}
\label{tab:versions-frontend}
\end{table}

\section{Hạ tầng chạy trên cụm}

Bảng sau ghi các image tag được khai báo cho cụm k3s. Đây là tag cấu hình, không nhất thiết thể hiện đầy đủ số hiệu bản vá của phần mềm bên trong image. PostgreSQL và Kafka ở môi trường này dùng dịch vụ được quản lý bên ngoài cụm, không có image tag trong cấu hình triển khai.

\begin{table}[H]
\centering
\small
\begin{tabularx}{\textwidth}{|p{3.4cm}|p{3.2cm}|X|}
\hline
Thành phần & Image tag & Nguồn / phạm vi \\ \hline
Keycloak & \textnormal{26.0} & Image tag khai báo trong cấu hình triển khai; không ghim số hiệu bản vá. \\ \hline
RustFS & \textnormal{1.0.5} & Image tag trong cấu hình triển khai tại \textnormal{deploy/rustfs}. \\ \hline
Docker-in-Docker & \textnormal{27-dind} & Image tag dùng cho executor trong Helm chart. \\ \hline
Kafka UI & \textnormal{v1.5.0} & Image tag trong cấu hình triển khai tại \textnormal{deploy/kafka-ui}. \\ \hline
\end{tabularx}
\caption{Phiên bản image của hạ tầng chạy trên cụm k3s.}
\label{tab:versions-infrastructure}
\end{table}

\section{Công cụ cài ngoài kho mã nguồn}

Node.js, k3s, Helm, Cloudflare Tunnel và kubectl không được khai báo phiên bản trong kho mã nguồn mà cài từ gói của hệ điều hành. Bảng~\ref{tab:versions-tools} ghi phiên bản đang cài trên máy tại thời điểm viết báo cáo và lệnh dùng để xác minh lại.

\begin{table}[H]
\centering
\small
\begin{tabularx}{\textwidth}{|p{3.3cm}|p{2.3cm}|>{\raggedright\arraybackslash}X|}
\hline
Công cụ & Phiên bản & Nguồn xác minh \\ \hline
Node.js & 26.13.1 & Gói \textnormal{nodejs} trên Arch Linux, dùng để build giao diện; kiểm tra bằng \textnormal{node -v}. \\ \hline
npm & 12.2.2 & Đi kèm gói Node.js; kiểm tra bằng \textnormal{npm -v}. \\ \hline
k3s & v1.37.2+k3s1 & Gói \textnormal{k3s-bin} từ AUR trên máy chủ cụm; kiểm tra bằng \textnormal{k3s --version}. \\ \hline
kubectl & v1.37.2 & Cùng phiên bản với cụm; kiểm tra bằng \textnormal{kubectl version --client}. \\ \hline
Helm & v4.3.0 & CLI dùng để cài các thành phần hạ tầng; kiểm tra bằng \textnormal{helm version --short}. \\ \hline
Cloudflare Tunnel & 2026.11.0 & Tiến trình \textnormal{cloudflared} chạy cùng cụm; kiểm tra bằng \textnormal{cloudflared --version}. \\ \hline
Argo CD & Không ghim & Cài từ manifest nhánh \textnormal{stable} trong \textnormal{deploy/install-argocd.sh}; số hiệu thực tế đọc từ image của Deployment \textnormal{argocd-server} trong namespace \textnormal{argocd}. \\ \hline
\end{tabularx}
\caption{Phiên bản công cụ cài ngoài kho mã nguồn tại thời điểm viết báo cáo.}
\label{tab:versions-tools}
\end{table}

Các số hiệu trên mô tả trạng thái máy khi viết báo cáo, không phải ràng buộc của dự án: cài lại từ gói mới nhất có thể cho phiên bản khác. Muốn tái lập đúng môi trường đã đo, cần ghim phiên bản bằng công cụ tương ứng, ví dụ Corepack hoặc asdf cho Node.js, biến \textnormal{INSTALL\_K3S\_VERSION} cho k3s, gói Helm theo phiên bản cố định, và chart hoặc manifest có image tag cụ thể cho Argo CD và Cloudflare Tunnel.

\end{document}
